% Additive manuscript excerpt. See integration/README.md for placement.
\section{All-rational formal classification of the Herglotz companion}
\label{hjr:sec:classification}

The conductor-doubling continuation~\cite{hergj:continuation} extends the
classification of $J(2/q)$ to every rational argument. The terminology
below refers to the rationalized five-term calculus, not to the kernel of
numerical period evaluation.

\begin{theorem}[All-rational formal classification]
\label{hjr:thm:classification}
Let $p,q$ be coprime positive integers. If they have different parity,
write $e$ for the even entry and $o$ for the odd entry. Then formal
rational-dilogarithmic and formal logarithmic reducibility of $J(p/q)$
coincide, and hold exactly when
\[
 e^2\equiv\pm1\pmod o,\qquad o^2\equiv1\pmod{2e}.
\]
If both entries are odd, formal rational-dilogarithmic reducibility holds
exactly when $p,q\in\{1,3,5\}$, and formal logarithmic reducibility holds
only when $p=q=1$. Thus the six non-logarithmic formal rational-dilogarithmic
exceptions are $1/3,3,1/5,5,3/5,5/3$.
\end{theorem}

The two new structural statements are an odd-conductor identity
\[
 \beta_{2m}(a)=3\beta_m(a)-\beta_m(2a)-\beta_m(a/2)
 \quad(m\text{ odd},\ (a,2m)=1)
\]
and even-conductor rigidity: for even $m$, a symbol $\beta_{2m}(a)$
can belong to the exterior square from $\Q(\zeta_m)$ only when it is
zero. The latter follows by restricting the positive Gram-matrix model
to the new character layer. Norm separation and two finite-group
difference tests then prove the classification; complete proofs are in
\cite[Sections 2--4]{hergj:continuation}.

\begin{theorem}[Recurrences and height counts]
\label{hjr:thm:recurrences}
Apart from $J(1)$, all unordered logarithmic pairs are adjacent terms in
$w_0=1$, $w_1=c$, $w_{j+1}=c w_j\pm w_{j-1}$ with even $c\ge2$.
The signs share only the seeds $(1,c)$. If $N(B)$ counts unordered pairs
of height at most $B$, then, for every fixed $R\ge2$,
\[
 N(B)=\frac32 B+\sum_{r=2}^{R}B^{1/r}
       +O_R\bigl(B^{1/(R+1)}\log B\bigr).
\]
\end{theorem}

\begin{theorem}[Full analytic evaluations]
\label{hjr:thm:evaluations}
Put $S(m)=\sum_{k=1}^{m-1}\log^2(2\sin(\pi k/m))$, with $S(1)=0$.
For every $n\ge1$,
\[
\begin{split}
 J(n/(n+1))={}&\frac12\log^22
 +\frac12\bigl(S(n)-S(n+1)+S(2n+2)-S(2n)\bigr)\\
 &-\frac{\pi^2(n^2-n-1)}{24n(n+1)}.
\end{split}
\]
With $\phi=(1+\sqrt5)/2$,
\[
\begin{split}
 J(3/5)={}&\Li_2(1/3)-\frac12\Li_2(1/5)
 +\frac12\log^22+\frac12\log^23-\frac14\log^25\\
 &+\log^2\phi-\frac{7\pi^2}{180}.
\end{split}
\]
All logarithms here are real logarithms of positive algebraic numbers.
\end{theorem}

The proofs use the normalized determinant-$2$ Hecke equation and Rogers
five-term substitutions within $(0,1)$, retaining all constants
\cite[Section 6]{hergj:continuation}. In particular,
\[
 J(1/3)-J(1/5)+J(3/5)
 =\frac12\log^22+2\log^2\phi-\frac{2\pi^2}{45}.
\]

\begin{remark}[A preprint finite-field assertion]
\label{hjr:rem:audit}
The finite-field identity
$\beta_p(n)=\beta_p(n+1)+\beta_p(n/(n+1))$ printed in
Section 7.2 of the reviewed Radchenko--Zagier preprint fails at $p=7,n=2$:
its defect is $3\beta_7(2)\ne0$. A strict two-embedding logarithmic
determinant proves nonvanishing. This does not affect the direct rational
boundary formula used in this chapter; see
\cite[Section 7]{hergj:continuation} for the exact witness and version scope.
\end{remark}
